% Generated by the blueprint tool (blueprint-pdf/1) from a blueprint snapshot; edit the
% Markdown sources, not this file.  Compile with `latexmk -pdf` or by running
% `pdflatex` twice (the second run settles the contents and the page references).
% Images, if any, are read from assets/ next to this file.
\documentclass[11pt]{article}
\usepackage[T1]{fontenc}
\usepackage[utf8]{inputenc}
\usepackage{lmodern}
\usepackage{textcomp}
\usepackage[a4paper,margin=2.5cm,headheight=14pt]{geometry}
\usepackage{amsmath,amssymb,amsthm}
\usepackage{amscd}  % KaTeX's CD diagrams; part of amsmath, so always there
\usepackage{mathtools}
\usepackage{bm}
\usepackage{graphicx}
\usepackage{xcolor}
\usepackage{array,booktabs,tabularx}
\usepackage{fancyvrb}
\usepackage[normalem]{ulem}
\usepackage{microtype}
% Optional, and only with their fonts: a small TeX installation has the
% packages without the fonts, which fails at the first formula.
\IfFileExists{mathrsfs.sty}{\IfFileExists{rsfs10.tfm}{\usepackage{mathrsfs}}{}}{}
\IfFileExists{stmaryrd.sty}{\IfFileExists{stmary10.tfm}{\usepackage{stmaryrd}}{}}{}
\IfFileExists{cancel.sty}{\usepackage{cancel}}{}
\usepackage[hyphens]{url}
\usepackage[pdfusetitle,bookmarksnumbered=false]{hyperref}
\usepackage{bookmark}
\definecolor{bplink}{HTML}{1F4E9C}
\definecolor{bpmuted}{HTML}{5F676E}
\hypersetup{colorlinks=true,linkcolor=bplink,urlcolor=bplink,citecolor=bplink}
\setlength{\parindent}{0pt}
\setlength{\parskip}{0.55em plus 0.2em minus 0.1em}
% Lists as tight as the page's.
\makeatletter
\def\@listi{\leftmargin\leftmargini\parsep 0pt\topsep 0.3em plus 0.1em\itemsep 0.1em plus 0.1em}
\let\@listI\@listi
\def\@listii{\leftmargin\leftmarginii\labelwidth\leftmarginii\advance\labelwidth-\labelsep\parsep 0pt\topsep 0.1em\itemsep 0.1em}
% An ordered list that does not start at 1.
\newcommand{\bpstartat}[1]{\setcounter{\@enumctr}{\numexpr#1-1\relax}}
\makeatother
\setcounter{tocdepth}{3}
\emergencystretch=3em
\allowdisplaybreaks

% A numbered statement, set as a paper sets one: "Theorem 2.1 (Title)."
% #1 is the lead ("Theorem~2.1"), #2 what follows it in roman.
\newenvironment{bpstatement}[2]{\par\addvspace{0.6\baselineskip}\noindent\textbf{#1}#2\textbf{.}\enspace\ignorespaces}
  {\par\addvspace{0.4\baselineskip}}
\newenvironment{bpstatementit}[2]{\begin{bpstatement}{#1}{#2}\itshape\ignorespaces}{\end{bpstatement}}
% A step: a relation with prose of its own, read under the object it starts from.
\newenvironment{bpstep}[1]{\par\addvspace{0.3\baselineskip}\begin{list}{}{\leftmargin=1.6em\rightmargin=0pt\topsep=0pt}\item\relax
  {\small\color{bpmuted}#1}\par\nopagebreak}{\end{list}\par}
% Notes about where an object is written out.
\newenvironment{bpnote}{\par\small\color{bpmuted}\itshape}{\par}
\newcommand{\bpomitted}[1]{\par{\small\color{bpmuted}\itshape #1\par}}
\newcommand{\bpbodyheading}[1]{\par\addvspace{0.4\baselineskip}\noindent\textbf{#1}\par\nopagebreak}
\newcommand{\bprule}{\par\medskip\noindent\hfil\rule{0.3\linewidth}{0.4pt}\hfil\par\medskip}
\newcommand{\bpboundary}[1]{\par\nopagebreak{\small\color{bpmuted}#1\par}\nopagebreak}
% A reference to something this export leaves out, and one that resolves nowhere.
\newcommand{\bpout}[1]{#1\textsuperscript{\hyperlink{bp-outside}{\ensuremath{\ast}}}}
\newcommand{\bpbroken}[1]{\textcolor{red!70!black}{[#1?]}}
\newcommand{\bpslug}[1]{\texttt{#1}}
% A bibliography entry: label, anchor, entry; the label hangs in the margin.
\newcommand{\bpbibitem}[3]{\par\setbox0\hbox{[#1]\enspace}\ifdim\wd0<5em\setbox0\hbox to 5em{[#1]\hfil}\fi
  \hangindent=\wd0 \hangafter=1 \noindent\box0\phantomsection\label{#2}#3\par\smallskip}
% A character the export knows no way to set.
\newcommand{\bpglyph}[1]{\fbox{\scriptsize #1}}
% Images, at their own size unless they are wider than the line.
\newsavebox{\bpimagebox}
\newcommand{\bpimage}[1]{\sbox{\bpimagebox}{\includegraphics{#1}}%
  \ifdim\wd\bpimagebox>\linewidth\resizebox{\linewidth}{!}{\usebox{\bpimagebox}}\else\usebox{\bpimagebox}\fi}
\newcommand{\bpmissingimage}[1]{\fbox{\parbox{0.8\linewidth}{\centering\small\color{bpmuted}Image not included: #1}}}
\DefineVerbatimEnvironment{bpverbatim}{Verbatim}{fontsize=\small,xleftmargin=1em}
\renewcommand{\qedsymbol}{\ensuremath{\blacksquare}}
\hypersetup{pdftitle={Weil Conjectures — B Abelian varieties},pdfsubject={Appendix B Abelian varieties; all levels},pdfcreator={blueprint}}
% KaTeX commands the website knows
\providecommand{\R}{\mathbb{R}}\providecommand{\N}{\mathbb{N}}\providecommand{\Z}{\mathbb{Z}}
\providecommand{\Q}{\mathbb{Q}}\providecommand{\C}{\mathbb{C}}
\providecommand{\Reals}{\mathbb{R}}\providecommand{\reals}{\mathbb{R}}\providecommand{\Complex}{\mathbb{C}}
\providecommand{\cnums}{\mathbb{C}}\providecommand{\natnums}{\mathbb{N}}\providecommand{\Bbb}[1]{\mathbb{#1}}
\providecommand{\bold}[1]{\mathbf{#1}}\providecommand{\frak}[1]{\mathfrak{#1}}\providecommand{\mathbbm}[1]{\mathbb{#1}}
\providecommand{\lang}{\langle}\providecommand{\rang}{\rangle}\providecommand{\bra}[1]{\langle #1\rvert}
\providecommand{\ket}[1]{\lvert #1\rangle}\providecommand{\braket}[1]{\langle #1\rangle}
\providecommand{\set}[1]{\{#1\}}\providecommand{\Set}[1]{\left\{#1\right\}}
\providecommand{\infin}{\infty}\providecommand{\isin}{\in}\providecommand{\plusmn}{\pm}\providecommand{\sub}{\subset}
\providecommand{\sube}{\subseteq}\providecommand{\supe}{\supseteq}\providecommand{\exist}{\exists}
\providecommand{\bull}{\bullet}\providecommand{\larr}{\leftarrow}\providecommand{\rarr}{\rightarrow}
\providecommand{\lrarr}{\leftrightarrow}\providecommand{\Larr}{\Leftarrow}\providecommand{\Rarr}{\Rightarrow}
\providecommand{\Lrarr}{\Leftrightarrow}\providecommand{\uarr}{\uparrow}\providecommand{\darr}{\downarrow}
\providecommand{\harr}{\leftrightarrow}\providecommand{\lArr}{\Leftarrow}\providecommand{\rArr}{\Rightarrow}
\providecommand{\hArr}{\Leftrightarrow}\providecommand{\alef}{\aleph}\providecommand{\alefsym}{\aleph}
\providecommand{\real}{\Re}\providecommand{\image}{\Im}\providecommand{\weierp}{\wp}
\providecommand{\Alpha}{\mathrm{A}}\providecommand{\Beta}{\mathrm{B}}\providecommand{\Epsilon}{\mathrm{E}}
\providecommand{\Zeta}{\mathrm{Z}}\providecommand{\Eta}{\mathrm{H}}\providecommand{\Iota}{\mathrm{I}}
\providecommand{\Kappa}{\mathrm{K}}\providecommand{\Mu}{\mathrm{M}}\providecommand{\Nu}{\mathrm{N}}
\providecommand{\Omicron}{\mathrm{O}}\providecommand{\Rho}{\mathrm{P}}\providecommand{\Tau}{\mathrm{T}}
\providecommand{\Chi}{\mathrm{X}}\providecommand{\omicron}{o}\providecommand{\gt}{>}\providecommand{\lt}{<}
\providecommand{\And}{\&}\providecommand{\hdashline}{\hline}\providecommand{\operatornamewithlimits}[1]{\operatorname*{#1}}
\providecommand{\llbracket}{[\![}\providecommand{\rrbracket}{]\!]}\providecommand{\cancel}[1]{#1}
\providecommand{\mathscr}{\mathcal}
% the characters beyond ASCII this document uses
\DeclareUnicodeCharacter{2162}{\textup{III}}
% the project's macros ([katex.macros] in blueprint.toml)
\providecommand{\A}{}\renewcommand{\A}{\mathbf A}
\providecommand{\AEn}{}\renewcommand{\AEn}{A_{E,n}}
\providecommand{\Aaf}{}\renewcommand{\Aaf}{\mathbf A}
\providecommand{\Ab}{}\renewcommand{\Ab}{\mathsf{Ab}}
\providecommand{\AffSch}{}\renewcommand{\AffSch}{\operatorname{AffSch}}
\providecommand{\Aone}{}\renewcommand{\Aone}{\mathbf A^1}
\providecommand{\Bl}{}\renewcommand{\Bl}{\operatorname{Bl}}
\providecommand{\Br}{}\renewcommand{\Br}{\mathrm{Br}}
\providecommand{\C}{}\renewcommand{\C}{\mathbb C}
\providecommand{\Cech}{}\renewcommand{\Cech}{\check{H}}
\providecommand{\Coext}{}\renewcommand{\Coext}{\operatorname{Coext}}
\providecommand{\Cone}{}\renewcommand{\Cone}{\operatorname{Cone}}
\providecommand{\Cons}{}\renewcommand{\Cons}{\operatorname{Cons}}
\providecommand{\Cref}{}\renewcommand{\Cref}[1]{\text{#1}}
\providecommand{\D}{}\renewcommand{\D}{D}
\providecommand{\DLAFT}{}\renewcommand{\DLAFT}{\operatorname{FT}}
\providecommand{\DLAFr}{}\renewcommand{\DLAFr}{\operatorname{Fr}}
\providecommand{\DLAQ}{}\renewcommand{\DLAQ}{{\overline{\mathbf Q}_\ell}}
\providecommand{\DLRQ}{}\renewcommand{\DLRQ}{{\overline{\mathbf Q}_\ell}}
\providecommand{\Dbc}{}\renewcommand{\Dbc}{D^b_c}
\providecommand{\Dc}{}\renewcommand{\Dc}{D_{\mathrm{cons}}}
\providecommand{\Dctf}{}\renewcommand{\Dctf}{D^b_{ctf}}
\providecommand{\Dic}{}\renewcommand{\Dic}{D_{\IndCons}}
\providecommand{\Dplus}{}\renewcommand{\Dplus}{D^{+}}
\providecommand{\End}{}\renewcommand{\End}{\operatorname{End}}
\providecommand{\Et}{}\renewcommand{\Et}{\mathrm{\acute{e}t}}
\providecommand{\Ex}{}\renewcommand{\Ex}{\mathrm{Ex}}
\providecommand{\Ext}{}\renewcommand{\Ext}{\operatorname{Ext}}
\providecommand{\F}{}\renewcommand{\F}{\mathbb F}
\providecommand{\FT}{}\renewcommand{\FT}{\operatorname{FT}}
\providecommand{\Fbar}{}\renewcommand{\Fbar}{{\overline{\mathbf F}_q}}
\providecommand{\Fell}{}\renewcommand{\Fell}{\mathbf F_\ell}
\providecommand{\Fib}{}\renewcommand{\Fib}{\operatorname{Fib}}
\providecommand{\Fix}{}\renewcommand{\Fix}{\operatorname{Fix}}
\providecommand{\Fp}{}\renewcommand{\Fp}{{\mathbf F_{p}}}
\providecommand{\Fq}{}\renewcommand{\Fq}{{\mathbf F_q}}
\providecommand{\Fqn}{}\renewcommand{\Fqn}{{\mathbf F_{q^{n}}}}
\providecommand{\Frob}{}\renewcommand{\Frob}{\operatorname{Frob}}
\providecommand{\GL}{}\renewcommand{\GL}{\operatorname{GL}}
\providecommand{\Gal}{}\renewcommand{\Gal}{\operatorname{Gal}}
\providecommand{\Hens}{}\renewcommand{\Hens}{\operatorname{Hens}}
\providecommand{\Hilb}{}\renewcommand{\Hilb}{\operatorname{Hilb}}
\providecommand{\Hom}{}\renewcommand{\Hom}{\operatorname{Hom}}
\providecommand{\IndCons}{}\renewcommand{\IndCons}{\operatorname{IndCons}}
\providecommand{\Kar}{}\renewcommand{\Kar}{\operatorname{Kar}}
\providecommand{\LT}{}\renewcommand{\LT}{\operatorname{LT}}
\providecommand{\Lie}{}\renewcommand{\Lie}{\operatorname{Lie}}
\providecommand{\Loc}{}\renewcommand{\Loc}{\operatorname{Loc}}
\providecommand{\Mod}{}\renewcommand{\Mod}{\mathsf{Mod}}
\providecommand{\N}{}\renewcommand{\N}{\mathbb{N}}
\providecommand{\Needspace}{}\renewcommand{\Needspace}{\par}
\providecommand{\O}{}\renewcommand{\O}{{\mathcal{O}}}
\providecommand{\Ob}{}\renewcommand{\Ob}{\mathrm{Ob}}
\providecommand{\Oe}{}\renewcommand{\Oe}{\mathcal O_E}
\providecommand{\PP}{}\renewcommand{\PP}{\mathbf{P}}
\providecommand{\Perf}{}\renewcommand{\Perf}{\operatorname{Perf}}
\providecommand{\Pic}{}\renewcommand{\Pic}{\mathrm{Pic}}
\providecommand{\Pm}{}\renewcommand{\Pm}{\mathbf P^{m}}
\providecommand{\Pone}{}\renewcommand{\Pone}{\mathbf P^{1}}
\providecommand{\Pqn}{}\renewcommand{\Pqn}{\mathbf P^{n}}
\providecommand{\ProFin}{}\renewcommand{\ProFin}{\operatorname{ProFin}}
\providecommand{\Proet}{}\renewcommand{\Proet}{\mathrm{pro\acute{e}t}}
\providecommand{\Q}{}\renewcommand{\Q}{\mathbb Q}
\providecommand{\Qb}{}\renewcommand{\Qb}{\overline{\mathbf Q}_{\ell}}
\providecommand{\Qbar}{}\renewcommand{\Qbar}{\overline{\mathbf Q}}
\providecommand{\QbarX}{}\renewcommand{\QbarX}{\overline{\mathbf Q}_{\ell,X}}
\providecommand{\Qbarl}{}\renewcommand{\Qbarl}{{\overline{\mathbf Q}_{\ell}}}
\providecommand{\Qell}{}\renewcommand{\Qell}{{\mathbf{Q}_{\ell}}}
\providecommand{\Ql}{}\renewcommand{\Ql}{\mathbf Q_{\ell}}
\providecommand{\Qlb}{}\renewcommand{\Qlb}{{\overline{\mathbf Q}_\ell}}
\providecommand{\Qlbar}{}\renewcommand{\Qlbar}{\overline{\mathbf Q}_\ell}
\providecommand{\R}{}\renewcommand{\R}{\mathbb{R}}
\providecommand{\RG}{}\renewcommand{\RG}{\operatorname{\mathrm{R}\Gamma}}
\providecommand{\RGamma}{}\renewcommand{\RGamma}{R\Gamma}
\providecommand{\RGammac}{}\renewcommand{\RGammac}{R\Gamma_c}
\providecommand{\RHom}{}\renewcommand{\RHom}{R\mathcal{H}om}
\providecommand{\RHomX}{}\renewcommand{\RHomX}{R\mathcal H om_X}
\providecommand{\Rat}{}\renewcommand{\Rat}{\operatorname{IsRat}}
\providecommand{\Res}{}\renewcommand{\Res}{\operatorname{Res}}
\providecommand{\Rlim}{}\renewcommand{\Rlim}{R\!\varprojlim}
\providecommand{\SIm}{}\renewcommand{\SIm}{\operatorname{SIm}}
\providecommand{\Sch}{}\renewcommand{\Sch}{\mathsf{Sch}}
\providecommand{\Sh}{}\renewcommand{\Sh}{\operatorname{Sh}}
\providecommand{\Span}{}\renewcommand{\Span}{\operatorname{span}}
\providecommand{\Spec}{}\renewcommand{\Spec}{\operatorname{Spec}}
\providecommand{\Tr}{}\renewcommand{\Tr}{\operatorname{Tr}}
\providecommand{\Trsm}{}\renewcommand{\Trsm}{\operatorname{Tr}^{\mathrm{sm}}}
\providecommand{\Wt}{}\renewcommand{\Wt}{\operatorname{Wt}}
\providecommand{\Z}{}\renewcommand{\Z}{\mathbb{Z}}
\providecommand{\Zar}{}\renewcommand{\Zar}{\mathrm{zar}}
\providecommand{\Zell}{}\renewcommand{\Zell}{{\mathbf{Z}_{\ell}}}
\providecommand{\Zl}{}\renewcommand{\Zl}{{\mathbf Z_\ell}}
\providecommand{\affet}{}\renewcommand{\affet}[1]{{#1}^{\mathrm{aff}}_{\mathrm{ét}}}
\providecommand{\affproet}{}\renewcommand{\affproet}[1]{{#1}^{\mathrm{aff}}_{\mathrm{proét}}}
\providecommand{\boxt}{}\renewcommand{\boxt}{\boxtimes_A^{\mathrm L}}
\providecommand{\cF}{}\renewcommand{\cF}{\mathcal{F}}
\providecommand{\cHom}{}\renewcommand{\cHom}{R\mathcal{H}\!om}
\providecommand{\cI}{}\renewcommand{\cI}{\mathcal{I}}
\providecommand{\cL}{}\renewcommand{\cL}{\mathcal L}
\providecommand{\cM}{}\renewcommand{\cM}{\mathcal M}
\providecommand{\cN}{}\renewcommand{\cN}{\mathcal N}
\providecommand{\cO}{}\renewcommand{\cO}{\mathcal O}
\providecommand{\cT}{}\renewcommand{\cT}{\mathcal T}
\providecommand{\checked}{}\renewcommand{\checked}{\color{red}\aftergroup\restorecolor}
\providecommand{\coeff}{}\renewcommand{\coeff}{\operatorname{coeff}}
\providecommand{\coker}{}\renewcommand{\coker}{\operatorname{coker}}
\providecommand{\colim}{}\renewcommand{\colim}{\operatorname{colim}}
\providecommand{\constructedby}{}\renewcommand{\constructedby}[1]{\par\smallskip\noindent\textcolor{purple!60!black}{\textsf{Constructed by:} }
  \footnotesize\sloppy\nolinkurl{#1}\par}
\providecommand{\deepresult}{}\renewcommand{\deepresult}[1]{\par\smallskip\noindent\textcolor{orange!45!black}{
  \textsf{Blueprint role: named deep theorem node; see the cluster proof architecture.}}
  \par\noindent\textcolor{orange!45!black}{\textsf{Mathematical source:} }#1\par}
\providecommand{\dt}{}\renewcommand{\dt}{\mathbin{\otimes^{\mathbf L}}}
\providecommand{\dw}{}\renewcommand{\dw}{\operatorname{dw}}
\providecommand{\et}{}\renewcommand{\et}[1]{{#1}_{\mathrm{ét}}}
\providecommand{\ev}{}\renewcommand{\ev}{\operatorname{ev}}
\providecommand{\fib}{}\renewcommand{\fib}{\operatorname{fib}}
\providecommand{\formalizationnote}{}\renewcommand{\formalizationnote}[1]{\begin{quote}\small\raggedright\textbf{Formalization note.} #1\end{quote}}
\providecommand{\fpqc}{}\renewcommand{\fpqc}{X_{\mathrm{fpqc}}}
\providecommand{\gammacont}{}\renewcommand{\gammacont}{\Gamma_{\mathrm{cont}}}
\providecommand{\gm}{}\renewcommand{\gm}{\mathbf G_m}
\providecommand{\hocolim}{}\renewcommand{\hocolim}{\operatorname{hocolim}}
\providecommand{\holim}{}\renewcommand{\holim}{\operatorname{holim}}
\providecommand{\hrB}{}\renewcommand{\hrB}{\@ifnextchar]{\hrEnd}{\hrStep}}
\providecommand{\hrEnd}{}\renewcommand{\hrEnd}{\@firstoftwo{}}
\providecommand{\hrStep}{}\renewcommand{\hrStep}{\@firstoftwo{\hrB}}
\providecommand{\id}{}\renewcommand{\id}{\operatorname{id}}
\providecommand{\kunFl}{}\renewcommand{\kunFl}{\mathbf F_\ell}
\providecommand{\kunZn}{}\renewcommand{\kunZn}{\mathbf Z/n\mathbf Z}
\providecommand{\kuncd}{}\renewcommand{\kuncd}{\operatorname{cd}}
\providecommand{\kundim}{}\renewcommand{\kundim}{\dim}
\providecommand{\kunotimes}{}\renewcommand{\kunotimes}[1]{\otimes^{\mathbf L}_{#1}}
\providecommand{\kunsupp}{}\renewcommand{\kunsupp}{\operatorname{Supp}}
\providecommand{\lotimes}{}\renewcommand{\lotimes}{\otimes^{\mathrm L}}
\providecommand{\m}{}\renewcommand{\m}{{\mathfrak{m}}}
\providecommand{\n}{}\renewcommand{\n}{{\mathfrak{n}}}
\providecommand{\nagataGap}{}\renewcommand{\nagataGap}[1]{\par\noindent\textbf{Formalization gap.} #1\par}
\providecommand{\openobligation}{}\renewcommand{\openobligation}[1]{\par\smallskip\noindent\textcolor{red!65!black}{\textsf{Open obligation:} }
  \footnotesize\sloppy\nolinkurl{#1}\par}
\providecommand{\p}{}\renewcommand{\p}{{\mathfrak{p}}}
\providecommand{\pr}{}\renewcommand{\pr}{\operatorname{pr}}
\providecommand{\preshv}{}\renewcommand{\preshv}[1]{\mathrm{PreShv}(#1)}
\providecommand{\pro}{}\renewcommand{\pro}[1]{\mathrm{Pro}(#1)}
\providecommand{\proet}{}\renewcommand{\proet}{X_{\mathrm{proét}}}
\providecommand{\q}{}\renewcommand{\q}{{\mathfrak{q}}}
\providecommand{\rad}{}\renewcommand{\rad}{\operatorname{rad}}
\providecommand{\rank}{}\renewcommand{\rank}{\operatorname{rank}}
\providecommand{\red}{}\renewcommand{\red}{\color{red}}
\providecommand{\restorecolor}{}\renewcommand{\restorecolor}{\color{black}}
\providecommand{\retiredobligation}{}\renewcommand{\retiredobligation}[1]{\par\smallskip\noindent\textcolor{red!75!black}{\textsf{Retired invalid route:} }
  \footnotesize\sloppy\nolinkurl{#1}\par}
\providecommand{\rk}{}\renewcommand{\rk}{\operatorname{rk}}
\providecommand{\sep}{}\renewcommand{\sep}{\textnormal{sep}}
\providecommand{\shv}{}\renewcommand{\shv}[1]{\mathrm{Shv}(#1)}
\providecommand{\spec}{}\renewcommand{\spec}[1]{\mathrm{Spec}(#1)}
\providecommand{\stacks}{}\renewcommand{\stacks}[1]{\href{https://stacks.math.columbia.edu/tag/#1}{\textnormal{Stacks, Tag~#1}}}
\providecommand{\stacksproject}{}\renewcommand{\stacksproject}[1]{{\cite[\href{https://stacks.math.columbia.edu/tag/#1}{Tag #1}]{stacks-project}}}
\providecommand{\statusnote}{}\renewcommand{\statusnote}[1]{\begin{quote}\small\raggedright\textbf{Status.} #1\end{quote}}
\providecommand{\tildeideal}{}\renewcommand{\tildeideal}[2]{\widetilde{(#1, #2)}}
\providecommand{\tildering}{}\renewcommand{\tildering}[3]{{#1}_{\widetilde{(#2, #3)}}}
\providecommand{\tr}{}\renewcommand{\tr}{\operatorname{tr}}
\providecommand{\wt}{}\renewcommand{\wt}{\operatorname{wt}}

\begin{document}
\pagestyle{plain}
\begin{center}
{\LARGE\bfseries Weil Conjectures\par}
\vspace{0.7em}
{\large Appendix B: Abelian varieties\par}
\vspace{0.5em}
{\small\color{bpmuted}Snapshot 2026-09-26, revision 8e4cb88e2a\par}
\end{center}
\vspace{0.5em}
{\small\tableofcontents}
\bigskip\hrule\bigskip

% object: chap-Abelian-varieties (blueprint/chap-Abelian-varieties/_section.md)
\section*{\texorpdfstring{B\quad Abelian varieties}{B Abelian varieties}}
\phantomsection\label{bp1}%
\addcontentsline{toc}{section}{\texorpdfstring{B\quad Abelian varieties}{B Abelian varieties}}
% text: prose-B-01
\phantomsection\label{bp2}%
Two steps of the main proof come down to the torsion points of a Jacobian. By the Kummer sequence, \(H^1(C,\mu_n)\) of a curve is the \(n\)-torsion of its Jacobian. The first step, the finiteness of \(H^1\) needed for proper base change, is therefore a count of torsion points. The second step is the finiteness of geometric monodromy in rank one, where purity begins. It needs Frobenius to fix no class in \(H^1\) of a curve over \(\mathbf F_q\) (Proposition \bpout{14.3}).

For a complex torus \(\mathbf C^g/\Lambda\) the count is trivial: the \(n\)-torsion is \(\tfrac1n\Lambda/\Lambda\simeq(\mathbf Z/n)^{2g}\). For an abelian variety \(A\) of dimension \(g\) over an arbitrary field \(k\), however, there is no lattice, and the torsion points must be counted algebraically. This is what we do: we prove that for \(n\) invertible in \(k\) the group scheme \(A[n]\) is finite étale of rank \(n^{2g}\) (Proposition \hyperref[bp34]{B.12}). It follows that the Tate module \(T_\ell A=\varprojlim_mA[\ell^m](\bar k)\) is free of rank \(2g\) over \(\mathbf Z_\ell\) (Proposition \hyperref[bp40]{B.14}). Moreover, over \(\mathbf F_q\), no power of Frobenius has a nonzero fixed vector on \(T_\ell A\) (Proposition \hyperref[bp43]{B.15}). The reason is short: a fixed vector would be a compatible system of \(\ell\)-power torsion points of unbounded order in the finite group \(A_0(\mathbf F_{q^m})\).

The count rests on one formula. If \(\mathcal N\) is a symmetric ample invertible sheaf on \(A\), then

\[
[n]^*\mathcal N\simeq\mathcal N^{\otimes n^2}
\]

(Proposition \hyperref[bp31]{B.11}). On an elliptic curve with \(\mathcal N=\mathcal O(e)\) this says that the divisor \([n]^*(e)\) is linearly equivalent to \(n^2(e)\). When \(n\) is prime to the characteristic, \([n]\) is étale and \([n]^*(e)\) is the reduced divisor \(E[n]\). Comparing degrees then gives \(\#E[n]=n^2\).

In general the formula is used twice. First, it forces \([n]\) to have finite fibres. Indeed, \(\mathcal N^{\otimes n^2}\) has positive degree on every complete curve in \(A\), while \([n]^*\mathcal N\) is trivial on a curve contracted by \([n]\). Second, we compute the Euler characteristics \(\chi(A,[n]^*\mathcal N^{\otimes m})\) in two ways, and comparing their leading coefficients gives \(\deg[n]=n^{2g}\).

The formula itself comes from the theorem of the cube (Theorem \hyperref[bp17]{B.6}). Following [Stacks, \href{https://stacks.math.columbia.edu/tag/0BEZ}{0BEZ}], we prove the theorem of the cube together with the rigidity lemma. As a by-product, from the rigidity lemma we deduce that \(A\) is commutative; see also [\hyperref[bp45]{Mil86a}].

In this appendix all group schemes are of finite type over a field \(k\) of arbitrary characteristic.

% object: sec-The-theorem-of-the-cube (blueprint/chap-Abelian-varieties/sec-The-theorem-of-the-cube/_section.md), content hidden
\subsection*{\texorpdfstring{B.1\quad The theorem of the cube}{B.1 The theorem of the cube}}
\phantomsection\label{bp3}%
\addcontentsline{toc}{subsection}{\texorpdfstring{B.1\quad The theorem of the cube}{B.1 The theorem of the cube}}
% text: prose-B-02
\phantomsection\label{bp4}%
The theorem of the cube says that an invertible sheaf \(\mathcal L\) on \(X\times Y\times Z\), with \(X,Y\) complete varieties and \(Z\) connected, that is trivial on the three slices through a point \((x,y,z)\) is trivial. To prove it, we view \(\mathcal L\) as a family of invertible sheaves on \(X\times Y\) parametrized by \(Z\). Following [Stacks, \href{https://stacks.math.columbia.edu/tag/0BEZ}{0BEZ}], we then consider the set of points of \(Z\) over which the sheaf is trivial. This set is closed, by semicontinuity. It is also open, by a computation of \(H^0\) and \(H^1\) near such a point. Since it contains \(z\) and \(Z\) is connected, it is all of \(Z\). The openness uses the triviality on the two slices \(x\times Y\times Z\) and \(X\times y\times Z\), and with only two factors it fails.

% object: example-cube-two-factors (blueprint/chap-Abelian-varieties/sec-The-theorem-of-the-cube/example-cube-two-factors/_section.md)
\phantomsection\label{bp5}%
\begin{bpstatement}{Example~B.1}{ (Two factors do not suffice)}
Let \(E\) be an elliptic curve over an algebraically closed field, with origin \(O\), and let \(\Delta\subset E\times E\) be the diagonal. On \(E\times E\) put

\[
\mathcal L=\mathcal O_{E\times E}\bigl(\Delta-E\times\{O\}-\{O\}\times E\bigr).
\]

Consider the slice \(E\times\{t\}\). It meets the diagonal in the point \(t\), and the divisor \(\{O\}\times E\) meets it in the point \(O\). The divisor \(E\times\{O\}\), on the other hand, is pulled back from the second factor, so it contributes the trivial sheaf for every \(t\). Hence the restriction is \(\mathcal O_E(t-O)\). By symmetry the restriction to \(\{t\}\times E\) is also \(\mathcal O_E(t-O)\). In particular \(\mathcal L\) is trivial on both slices through \((O,O)\).

However, \(\mathcal L\) is not trivial, because for \(t\neq O\) the sheaf \(\mathcal O_E(t-O)\) is nontrivial. Indeed, a rational function with divisor \(t-O\) would be a map \(E\to\mathbf P^1\) of degree one, which is impossible in genus one. The set of \(t\) over which the fibre is trivial is the point \(O\). It is closed, as Lemma \hyperref[bp6]{B.2} below predicts, but not open.
\end{bpstatement}

% object: lemma-triviality-fibres (blueprint/chap-Abelian-varieties/sec-The-theorem-of-the-cube/lemma-triviality-fibres/_section.md)
\phantomsection\label{bp6}%
\begin{bpstatementit}{Lemma~B.2}{ (Triviality on fibres)}
Let \(P\) be a proper geometrically integral \(k\)-scheme, let \(T\) be a \(k\)-scheme of finite type, let \(f:P\times T\to T\) be the projection, and let \(\mathcal L\) be an invertible sheaf on \(P\times T\). For \(t\in T\) write \(P_t=P\times_k\kappa(t)\) and \(\mathcal L_t=\mathcal L|_{P_t}\).

\begin{enumerate}
\item \(\mathcal L_t\) is trivial if and only if \(H^0(P_t,\mathcal L_t)\ne0\) and \(H^0(P_t,\mathcal L_t^{-1})\ne0\).
\item The set of \(t\in T\) such that \(\mathcal L_t\) is trivial is closed.
\end{enumerate}
\end{bpstatementit}

% proof: proof/lemma-triviality-fibres
\begin{proof}[{Proof}]\phantomsection\label{bp7}%
The proof uses only that \(P_t\) is proper and integral with \(H^0(P_t,\mathcal O_{P_t})=\kappa(t)\), which holds by [Stacks, \href{https://stacks.math.columbia.edu/tag/0BUG}{0BUG}] because \(P\) is proper and geometrically integral.

For (1), a trivial sheaf and its inverse both have the nonzero section \(1\), so only the converse needs an argument. On \(\mathbf P^1\), for instance, \(\mathcal O(d)\) and \(\mathcal O(-d)\) both have nonzero sections only when \(d=0\), and the general argument is the same. Let \(s\) and \(s'\) be nonzero sections of \(\mathcal L_t\) and \(\mathcal L_t^{-1}\). Since \(P_t\) is integral, the product \(ss'\) is a nonzero element of \(H^0(P_t,\mathcal O)=\kappa(t)\), hence a unit. Since \(ss'\) would vanish wherever \(s\) does, the section \(s\) vanishes nowhere, and it trivializes \(\mathcal L_t\).

For (2), note that by (1) the set in question is the intersection of the loci where \(\dim H^0(P_t,\mathcal L_t)\) and \(\dim H^0(P_t,\mathcal L_t^{-1})\) are positive. These loci are closed because the functions \(t\mapsto\dim H^0(P_t,\mathcal L_t^{\pm1})\) are upper semicontinuous [Stacks, \href{https://stacks.math.columbia.edu/tag/0BDN}{0BDN}], as \(\mathcal L\) is flat over \(T\) and \(f\) is proper.
\end{proof}

% text: prose-B-03
\phantomsection\label{bp8}%
Part (2) is the closed half of the theorem of the cube. If the sheaf is trivial over the members of an open cover of \(T\), then the evaluation map identifies it with a pullback from \(T\), since sections along the proper factor \(P\) are constant.

% object: corollary-triviality-over-open-cover-base (blueprint/chap-Abelian-varieties/sec-The-theorem-of-the-cube/corollary-triviality-over-open-cover-base/_section.md)
\phantomsection\label{bp9}%
\begin{bpstatementit}{Corollary~B.3}{ (Triviality over an open cover of the base)}
With the notation of Lemma \hyperref[bp6]{B.2}, suppose that \(\mathcal L\) is trivial on \(P\times U_\alpha\) for an open cover \((U_\alpha)\) of \(T\). Then \(f_*\mathcal L\) is invertible and the evaluation map \(f^*f_*\mathcal L\to\mathcal L\) is an isomorphism.
\end{bpstatementit}

% proof: proof/corollary-triviality-over-open-cover-base
\begin{proof}[{Proof}]\phantomsection\label{bp10}%
Both assertions are local on \(T\), so we may replace \(T\) by one of the \(U_\alpha\) and assume that \(\mathcal L\simeq\mathcal O_{P\times T}\). Since \(H^0(P,\mathcal O_P)=k\) [Stacks, \href{https://stacks.math.columbia.edu/tag/0BUG}{0BUG}], and \(T\) is flat over \(k\), flat base change gives \(f_*\mathcal O_{P\times T}=\mathcal O_T\). Hence the chosen trivialization identifies \(f_*\mathcal L\) with the invertible sheaf \(\mathcal O_T\). Under the same identification the evaluation map \(f^*f_*\mathcal L\to\mathcal L\) becomes the identity of \(\mathcal O_{P\times T}\).
\end{proof}

% text: prose-B-04
\phantomsection\label{bp11}%
The open half of the theorem is the substantial step. It uses the fact that \(H^1\) of the structure sheaf of a product is the sum of the \(H^1\) of the factors.

% object: lemma-h-1-structure-sheaf-product (blueprint/chap-Abelian-varieties/sec-The-theorem-of-the-cube/lemma-h-1-structure-sheaf-product/_section.md)
\phantomsection\label{bp12}%
\begin{bpstatementit}{Lemma~B.4}{ ($H^1$ of the structure sheaf of a product)}
Let \(K\) be a field, let \(P\) and \(Q\) be proper \(K\)-schemes with \(H^0(P,\mathcal O_P)=H^0(Q,\mathcal O_Q)=K\), and let \(p_0\in P(K)\) and \(q_0\in Q(K)\). Restriction to the slices \(P\times q_0\) and \(p_0\times Q\) is an isomorphism

\[
H^1(P\times Q,\mathcal O)\longrightarrow
H^1(P,\mathcal O)\oplus H^1(Q,\mathcal O).
\]
\end{bpstatementit}

% proof: proof/lemma-h-1-structure-sheaf-product
\begin{proof}[{Proof}]\phantomsection\label{bp13}%
By the Künneth formula [Stacks, \href{https://stacks.math.columbia.edu/tag/0BED}{0BED}], and since \(H^0\) of \(\mathcal O\) is \(K\) on both factors, the source is \(\operatorname{pr}_1^*H^1(P,\mathcal O)\oplus
\operatorname{pr}_2^*H^1(Q,\mathcal O)\). Restriction to \(P\times q_0\) is the identity on the first summand. It kills the second, because on the second summand it factors through \(H^1\) of the point \(q_0\). Restriction to \(p_0\times Q\) does the opposite.
\end{proof}

% object: lemma-triviality-near-fibre (blueprint/chap-Abelian-varieties/sec-The-theorem-of-the-cube/lemma-triviality-near-fibre/_section.md)
\phantomsection\label{bp14}%
\begin{bpstatementit}{Lemma~B.5}{ (Triviality near a fibre)}
Let \(P\) and \(Q\) be proper geometrically integral \(k\)-schemes with points \(p_0\in P(k)\) and \(q_0\in Q(k)\), let \(T\) be a \(k\)-scheme of finite type, and let \(\mathcal L\) be an invertible sheaf on \(P\times Q\times T\) whose restrictions to \(p_0\times Q\times T\) and to \(P\times q_0\times T\) are trivial. If \(\mathcal L\) is trivial on the fibre over a point \(t\in T\), then \(\mathcal L\) is trivial on \(P\times Q\times U\) for an open neighbourhood \(U\) of \(t\).
\end{bpstatementit}

% proof: proof/lemma-triviality-near-fibre
\begin{proof}[{Proof}]\phantomsection\label{bp15}%
We want a section of \(\mathcal L\) over \(X\times U\), where \(X=P\times Q\), that does not vanish on the fibre over \(t\). Its zero locus then misses a neighbourhood of that fibre. Such a section exists as soon as \(f_*\mathcal L\) is free of rank one near \(t\), with formation compatible with passage to the fibre.

To get this, we represent \(Rf_*\mathcal L\), and the pushforwards of the restrictions of \(\mathcal L\) to the two slices, by complexes of free modules whose differentials vanish at \(t\). By the Künneth formula, restriction from \(H^1\) of the fibre to \(H^1\) of the two slices is an isomorphism. This forces the first differential of \(Rf_*\mathcal L\) to vanish after shrinking \(T\). Then \(H^0\) of \(Rf_*\mathcal L\) is free of rank one, as we need.

\emph{Setup.} Since \(X_{\bar k}=P_{\bar k}\times Q_{\bar k}\) is a product of integral varieties, \(X\) is geometrically integral [Stacks, \href{https://stacks.math.columbia.edu/tag/05P3}{05P3}]. We may assume that \(T=\operatorname{Spec}R\) is affine, and we let \(\mathfrak p\) be the prime of \(t\). Below, “shrinking \(T\)” means replacing \(R\) by \(R_g\) for some \(g\notin\mathfrak p\); everything already constructed survives this. Consider the projections

\[
f:X\times T\to T,\qquad f_1:P\times T\to T,\qquad f_2:Q\times T\to T,
\]

and the restrictions \(\mathcal L_1\) and \(\mathcal L_2\) of \(\mathcal L\) to \(P\times q_0\times T\) and \(p_0\times Q\times T\), which are trivial. The complexes \(E=Rf_*\mathcal L\) and \(E_i=Rf_{i*}\mathcal L_i\) are perfect, and their formation commutes with base change [Stacks, \href{https://stacks.math.columbia.edu/tag/0B91}{0B91}]. Restriction of sections gives a map \(E\to E_1\oplus E_2\), also compatible with base change. Since \(T\) is affine, \(H^0(X\times T,\mathcal L)=H^0(E)\) and likewise for \(E_i\).

\emph{Minimal complexes.} A perfect complex over a local ring can be represented by a complex of free modules whose differentials vanish modulo the maximal ideal; the ranks are then the dimensions of the cohomology of the fibre. We use this to read off the ranks from the fibre over \(t\), where everything is trivial. To do so, put \(h^j=\dim H^j(X_t,\mathcal O)\) and \(h_i^j=\dim H^j\) of \(\mathcal O\) on the fibres \(P_t\) and \(Q_t\). Since \(\mathcal L\) is trivial on \(X_t\) and the \(\mathcal L_i\) are trivial, \(E\otimes^{\mathbf L}_R\kappa(t)\) and \(E_i\otimes^{\mathbf L}_R\kappa(t)\) have cohomology of dimensions \(h^j\) and \(h_i^j\), with \(h^0=h_i^0=1\). After shrinking \(T\) we may therefore represent \(E\) and \(E_i\) by complexes of finite free modules

\[
R\xrightarrow{d}R^{h^1}\to R^{h^2}\to\cdots,\qquad
R\xrightarrow{d_i}R^{h_i^1}\to R^{h_i^2}\to\cdots
\]

whose differentials vanish modulo \(\mathfrak p\) [Stacks, \href{https://stacks.math.columbia.edu/tag/0BCD}{0BCD}]. We represent the map \(E\to E_i\) by a map of complexes [Stacks, \href{https://stacks.math.columbia.edu/tag/064B}{064B}], with components \(a_i:R\to R\) in degree \(0\) and \(b_i:R^{h^1}\to R^{h_i^1}\) in degree \(1\). Since the differentials vanish modulo \(\mathfrak p\), the map \(b_i\otimes\kappa(t)\) is the restriction map on \(H^1\) of the fibre over \(t\).

We want to show that \(d=0\) after shrinking \(T\), since then \(H^0(E)=\ker d=R\) is free of rank one. The maps \(b_i\) will let us transfer this question to the \(d_i\), which are easier because \(\mathcal L_i\) is trivial.

\emph{The differentials \(d_i\) vanish.} Fix a trivialization \(\mathcal L_1\simeq\mathcal O\), and let \(u\in H^0(E_1)=\ker d_1\subset R\) be the section corresponding to \(1\). Its image in \(H^0(E_1\otimes^{\mathbf L}\kappa(t))=H^0(P_t,\mathcal O)\) is \(1\ne0\), so \(u\notin\mathfrak p\). Since \(d_1(u)=u\,d_1(1)=0\), after inverting \(u\) we get \(d_1=0\). The same argument gives \(d_2=0\).

\emph{The differential \(d\) vanishes.} This step uses the two slices together. If the map \((b_1,b_2)\) is an isomorphism, then \(d\) is determined by \(d_1\) and \(d_2\), and this condition can be checked on the fibre. The map \((b_1,b_2)\otimes\kappa(t)\) is restriction from \(H^1(P_t\times Q_t,\mathcal O)\) to the slices \(P_t\times q_0\) and \(p_0\times Q_t\). It is an isomorphism by Lemma \hyperref[bp12]{B.4}, and in particular \(h^1=h_1^1+h_2^1\). Hence the square matrix \((b_1,b_2)\) has determinant outside \(\mathfrak p\). After inverting this determinant, \((b_1,b_2)\) is an isomorphism. Since the maps form a map of complexes, we have \(b_i\circ d=d_i\circ a_i=0\), and therefore \(d=0\).

\emph{Conclusion.} A generator of \(H^0\) gives the trivialization. We have \(H^0(X\times T,\mathcal L)=H^0(E)=\ker d=R\), and its generator \(\sigma\) maps to a generator of \(H^0(E\otimes^{\mathbf L}\kappa(t))=H^0(X_t,\mathcal L_t)\). Since \(\mathcal L_t\simeq\mathcal O_{X_t}\) and \(H^0(X_t,\mathcal O)=\kappa(t)\), the section \(\sigma\) vanishes nowhere on \(X_t\). Its zero locus is closed in \(X\times T\) and disjoint from \(X_t\). Since \(f\) is proper, its image in \(T\) is a closed set not containing \(t\). On the complement \(U\) of this image, \(\sigma\) trivializes \(\mathcal L\).
\end{proof}

% text: prose-B-05
\phantomsection\label{bp16}%
With the closed and the open half in hand, it remains to glue the local trivializations.

% object: theorem-theorem-cube (blueprint/chap-Abelian-varieties/sec-The-theorem-of-the-cube/theorem-theorem-cube/_section.md)
\phantomsection\label{bp17}%
\begin{bpstatementit}{Theorem~B.6}{ (Theorem of the cube)}
Let \(X\) and \(Y\) be proper geometrically integral \(k\)-schemes with points \(x\in X(k)\) and \(y\in Y(k)\), let \(Z\) be a connected \(k\)-scheme of finite type with a point \(z\in Z(k)\), and let \(\mathcal L\) be an invertible sheaf on \(X\times Y\times Z\). If the restrictions of \(\mathcal L\) to \(x\times Y\times Z\), \(X\times y\times Z\) and \(X\times Y\times z\) are trivial, then \(\mathcal L\) is trivial.
\end{bpstatementit}

% proof: proof/theorem-theorem-cube
\begin{proof}[{Proof}]\phantomsection\label{bp18}%
The local trivializations are glued by Corollary \hyperref[bp9]{B.3}, which shows that \(\mathcal L\) is pulled back from \(Y\times Z\); the slice \(x\times Y\times Z\) then identifies the sheaf on \(Y\times Z\).

Let \(Z_0\subset Z\) be the set of points over which \(\mathcal L\) is trivial on the fibre of \(p:X\times Y\times Z\to Z\). It contains \(z\). It is closed by Lemma lemma-triviality-fibres applied to \(X\times Y\), which is proper and geometrically integral [Stacks, \href{https://stacks.math.columbia.edu/tag/05P3}{05P3}]. It is open by Lemma \hyperref[bp14]{B.5}. Since \(Z\) is connected, \(Z_0=Z\). Now Lemma \hyperref[bp14]{B.5} gives an open cover of \(Z\) over whose members \(\mathcal L\) is trivial.

Let \(\pi:X\times Y\times Z\to Y\times Z\) be the projection and \(\sigma\) its section through \(x\). We apply Corollary \hyperref[bp9]{B.3} with \(P=X\), \(T=Y\times Z\) and the open cover by the products of \(Y\) with these opens. It shows that \(\mathcal M=\pi_*\mathcal L\) is invertible and that \(\pi^*\mathcal M\to\mathcal L\) is an isomorphism. Hence \(\mathcal M\simeq\sigma^*\pi^*\mathcal M\simeq\sigma^*\mathcal L\). This is the restriction of \(\mathcal L\) to \(x\times Y\times Z\), which is trivial by hypothesis. We conclude that \(\mathcal L\simeq\pi^*\mathcal M\) is trivial.
\end{proof}

% text: prose-B-06
\phantomsection\label{bp19}%
The second tool, which we use to show that an abelian variety is commutative (Lemma \hyperref[bp26]{B.9}), is a rigidity statement. It says that a family of maps from a proper variety that is constant on one member is constant on all nearby members.

% object: lemma-rigidity (blueprint/chap-Abelian-varieties/sec-The-theorem-of-the-cube/lemma-rigidity/_section.md)
\phantomsection\label{bp20}%
\begin{bpstatementit}{Lemma~B.7}{ (Rigidity)}
Let \(X\) be a proper geometrically integral \(k\)-scheme, let \(Y\) be a \(k\)-scheme of finite type, let \(Z\) be a \(k\)-scheme, and let \(h:X\times Y\to Z\) be a morphism. If \(h\) maps \(X\times y_0\) to a single point for some \(y_0\in Y\), then there are an open neighbourhood \(V\) of \(y_0\) and a morphism \(g:V\to Z\) with \(h=g\circ\operatorname{pr}_2\) on \(X\times V\).
\end{bpstatementit}

% proof: proof/lemma-rigidity
\begin{proof}[{Proof}]\phantomsection\label{bp21}%
Near \(X\times y_0\) the map \(h\) lands in an affine open of \(Z\), and a map to an affine scheme is given by functions. Since \(X\) is proper and geometrically integral, functions on \(X\times V\) are functions on \(V\). Hence \(h\) factors through \(V\).

Let \(W\subset Z\) be an affine open containing \(h(X\times y_0)\). Since \(X\) is proper, the image in \(Y\) of the closed set \(h^{-1}(Z\setminus W)\) is closed, and its complement contains \(y_0\). Let \(V\) be an affine open neighbourhood of \(y_0\) in this complement. Then \(h\) maps \(X\times V\) into the affine scheme \(W\). It is therefore determined by the ring map \(\mathcal O(W)\to H^0(X\times V,\mathcal O)\) [Stacks, \href{https://stacks.math.columbia.edu/tag/01I1}{01I1}]. By flat base change [Stacks, \href{https://stacks.math.columbia.edu/tag/02KH}{02KH}] and \(H^0(X,\mathcal O_X)=k\) [Stacks, \href{https://stacks.math.columbia.edu/tag/0BUG}{0BUG}], we have \(H^0(X\times V,\mathcal O)=\mathcal O(V)\). The ring map is thus \(\operatorname{pr}_2^*\) composed with a ring map \(\mathcal O(W)\to\mathcal O(V)\), and this ring map defines \(g\).
\end{proof}


% object: sec-Multiplication-by-n (blueprint/chap-Abelian-varieties/sec-Multiplication-by-n/_section.md), content hidden
\subsection*{\texorpdfstring{B.2\quad Multiplication by $n$}{B.2 Multiplication by n}}
\phantomsection\label{bp22}%
\addcontentsline{toc}{subsection}{\texorpdfstring{B.2\quad Multiplication by $n$}{B.2 Multiplication by n}}
% text: prose-B-07
\phantomsection\label{bp23}%
Let \(\mathcal L\) be an invertible sheaf on an abelian variety \(A\). We pull \(\mathcal L\) back to \(A\times A\times A\) along the partial sums of the coordinates. There the theorem of the cube gives a second order recursion for \(n\mapsto[n]^*\mathcal L\), whose solution is

\[
[n]^*\mathcal L\simeq\mathcal L^{\otimes n(n+1)/2}\otimes([-1]^*\mathcal L)^{\otimes n(n-1)/2}
\]

(Proposition \hyperref[bp31]{B.11}). For symmetric \(\mathcal L\) the formula gives \(\mathcal L^{\otimes n^2}\). If \([-1]^*\mathcal L\simeq\mathcal L^{-1}\), as for an invertible sheaf of degree zero on an elliptic curve, it gives \(\mathcal L^{\otimes n}\).

We apply the formula to a symmetric ample sheaf (Lemma \hyperref[bp29]{B.10}). It shows that \([n]\) is finite and flat of degree \(n^{2g}\), and that \([n]\) is étale when \(n\) is invertible in \(k\) (Proposition \hyperref[bp34]{B.12}). The hypothesis on \(n\) cannot be dropped. For example, in characteristic \(p\) the differential of \([p]\) is zero, and as a consequence an elliptic curve has at most \(p\) points of order dividing \(p\) over \(\bar k\), although \(E[p]\) has rank \(p^2\).

% object: definition-abelian-variety (blueprint/chap-Abelian-varieties/sec-Multiplication-by-n/definition-abelian-variety/_section.md)
\phantomsection\label{bp24}%
\begin{bpstatement}{Definition~B.8}{ (Abelian variety)}
An \emph{abelian variety} over \(k\) is a smooth proper geometrically connected \(k\)-group scheme. If \(A\) is an abelian variety, \([n]\) denotes multiplication by the integer \(n\), and \(A[n]=\ker[n]\) for \(n>0\). An invertible sheaf \(\mathcal L\) on \(A\) is \emph{symmetric} if \([-1]^*\mathcal L\simeq\mathcal L\).
\end{bpstatement}

% text: prose-B-08
\phantomsection\label{bp25}%
We use three facts: commutativity, which follows from rigidity; a symmetric ample sheaf, which follows from projectivity; and the formula for \([n]^*\mathcal L\), which follows from the cube. Rigidity and the cube require \(A\) to be geometrically integral, which we check together with commutativity.

% object: lemma-abelian-variety-commutative (blueprint/chap-Abelian-varieties/sec-Multiplication-by-n/lemma-abelian-variety-commutative/_section.md)
\phantomsection\label{bp26}%
\begin{bpstatementit}{Lemma~B.9}{ (Commutativity)}
An abelian variety \(A\) over \(k\) is geometrically integral and commutative.
\end{bpstatementit}

% proof: proof/lemma-abelian-variety-commutative
\begin{proof}[{Proof}]\phantomsection\label{bp27}%
An abelian variety is geometrically reduced because it is smooth, and geometrically irreducible because the identity component of a group scheme over a field is geometrically irreducible [Stacks, \href{https://stacks.math.columbia.edu/tag/0B7R}{0B7R}] and \(A\) is connected.

For commutativity, consider \(h:A\times A\to A\), \(h(x,y)=xyx^{-1}y^{-1}\). It maps \(A\times\{e\}\) to \(e\). By Lemma \hyperref[bp20]{B.7}, applied with the first factor as the proper one, there are an open \(V\ni e\) and \(g:V\to A\) with \(h|_{A\times V}=g\circ\operatorname{pr}_2\). Composing with \((e,\operatorname{id}):V\to A\times V\) gives \(g=h(e,-)=e\), so \(h=e\) on \(A\times V\). To pass from \(A\times V\) to \(A\times A\), we compare \(h\) with the constant map \(e\). Since \(A\times A\) is smooth, hence reduced, and irreducible, the nonempty open \(A\times V\) is scheme-theoretically dense. Two maps to a separated scheme that agree on a scheme-theoretically dense open agree everywhere (Lemma \bpout{A.4}). Hence \(h=e\) [Stacks, \href{https://stacks.math.columbia.edu/tag/056D}{056D}, \href{https://stacks.math.columbia.edu/tag/01RH}{01RH}].
\end{proof}

% text: prose-B-09
\phantomsection\label{bp28}%
From now on we write the group law additively. Projectivity is the one structural fact that we quote.

% object: lemma-symmetric-ample-sheaf (blueprint/chap-Abelian-varieties/sec-Multiplication-by-n/lemma-symmetric-ample-sheaf/_section.md)
\phantomsection\label{bp29}%
\begin{bpstatementit}{Lemma~B.10}{ (A symmetric ample sheaf)}
An abelian variety \(A\) over \(k\) is projective, and it has a symmetric ample invertible sheaf.
\end{bpstatementit}

% proof: proof/lemma-symmetric-ample-sheaf
\begin{proof}[{Proof}]\phantomsection\label{bp30}%
Every group scheme of finite type over a field is quasi-projective [Stacks, \href{https://stacks.math.columbia.edu/tag/0BF7}{0BF7}]. A quasi-projective proper scheme is projective [Stacks, \href{https://stacks.math.columbia.edu/tag/0B45}{0B45}]. If \(\mathcal L\) is ample, then so is \([-1]^*\mathcal L\), since \([-1]\) is an automorphism, and the tensor product \(\mathcal N=\mathcal L\otimes[-1]^*\mathcal L\) of two ample sheaves is ample. It is symmetric because \([-1]\) is an involution, which interchanges the two factors.
\end{proof}

% object: proposition-pullback-multiplication-n (blueprint/chap-Abelian-varieties/sec-Multiplication-by-n/proposition-pullback-multiplication-n/_section.md)
\phantomsection\label{bp31}%
\begin{bpstatementit}{Proposition~B.11}{ (Pullback of invertible sheaves by \{$[n]$\})}
Let \(A\) be an abelian variety over \(k\) and let \(\mathcal L\) be an invertible sheaf on \(A\). For every integer \(n\geq0\),

\[
[n]^*\mathcal L\simeq
\mathcal L^{\otimes n(n+1)/2}\otimes
([-1]^*\mathcal L)^{\otimes n(n-1)/2}.
\]

In particular \([n]^*\mathcal N\simeq\mathcal N^{\otimes n^2}\) if \(\mathcal N\) is symmetric.
\end{bpstatementit}

% proof: proof/proposition-pullback-multiplication-n
\begin{proof}[{Proof}]\phantomsection\label{bp32}%
The idea is to produce one trivial sheaf \(\mathcal M\) on \(A\times A\times A\) by the theorem of the cube, and to pull it back along maps \(A\to A\times A\times A\) chosen so that the result is a relation between \([n+1]^*\mathcal L\), \([n]^*\mathcal L\) and \([n-1]^*\mathcal L\). The theorem of the cube applies because \(A\) is geometrically integral (Lemma \hyperref[bp26]{B.9}).

Consider on \(A\times A\times A\) the invertible sheaf

\[
\mathcal M=m_{123}^*\mathcal L\otimes m_{12}^*\mathcal L^{-1}\otimes
m_{13}^*\mathcal L^{-1}\otimes m_{23}^*\mathcal L^{-1}\otimes
p_1^*\mathcal L\otimes p_2^*\mathcal L\otimes p_3^*\mathcal L,
\]

where the \(p_i\) are the projections and \(m_{12}\), \(m_{13}\), \(m_{23}\), \(m_{123}\) are the sums of the indicated coordinates. To apply the theorem of the cube, we restrict \(\mathcal M\) to the three slices. On \(e\times A\times A\), with coordinates \((x_2,x_3)\), the seven factors restrict to \(\mathcal L\) pulled back along \(x_2+x_3\), \(x_2\), \(x_3\), \(x_2+x_3\), the constant map \(e\), \(x_2\) and \(x_3\), with exponents \(1,-1,-1,-1,1,1,1\). Of these, the pullback along the constant map is trivial, and the others cancel. The same holds on \(A\times e\times A\) and \(A\times A\times e\). Hence \(\mathcal M\) is trivial by Theorem \hyperref[bp17]{B.6}. Finally, pulling the trivial sheaf \(\mathcal M\) back along \(x\mapsto(x,x,[n-1]x)\) for \(n\ge0\) yields

\[
[n+1]^*\mathcal L\otimes\mathcal L^{\otimes2}
\otimes[n-1]^*\mathcal L
\simeq [2]^*\mathcal L\otimes([n]^*\mathcal L)^{\otimes2}.
\]

To see this, note that the composites of \(x\mapsto(x,x,[n-1]x)\) with \(m_{123}\), \(m_{12}\), \(m_{13}\), \(m_{23}\), \(p_1\), \(p_2\), \(p_3\) are \([n+1]\), \([2]\), \([n]\), \([n]\), \(\mathrm{id}\), \(\mathrm{id}\), \([n-1]\). The cases \(n=0\) and \(n=1\) of the formula are clear. For \(n=0\), since \([0]^*\mathcal L\) is trivial, the relation reads \(\mathcal L^{\otimes3}\otimes[-1]^*\mathcal L\simeq[2]^*\mathcal L\). This is the case \(n=2\) of the formula. In general, the relation expresses \([n+1]^*\mathcal L\) in terms of \([n]^*\mathcal L\) and \([n-1]^*\mathcal L\), and the exponents in the formula satisfy the same second order recursion. Induction on \(n\) therefore gives the formula [Stacks, \href{https://stacks.math.columbia.edu/tag/0BFF}{0BFF}]. In particular, for symmetric \(\mathcal N\) the two exponents add up to \(n^2\).
\end{proof}

% text: prose-B-10
\phantomsection\label{bp33}%
For a symmetric ample \(\mathcal N\), the degree of \([n]\) can be read off from the growth of \(\chi(A,\mathcal N^{\otimes m})\), since pulling back by \([n]\) multiplies \(m\) by \(n^2\) and \(\chi\) grows like \(m^g\).

% object: proposition-multiplication-n (blueprint/chap-Abelian-varieties/sec-Multiplication-by-n/proposition-multiplication-n/_section.md)
\phantomsection\label{bp34}%
\begin{bpstatementit}{Proposition~B.12}{ (Multiplication by $n$)}
Let \(A\) be an abelian variety of dimension \(g\) over \(k\). For every \(n>0\), the morphism \([n]\) is surjective and finite locally free of degree \(n^{2g}\), and its kernel \(A[n]\) has rank \(n^{2g}\) over \(k\). If \(n\) is invertible in \(k\), then \([n]\) and \(A[n]\to\operatorname{Spec}k\) are finite étale.
\end{bpstatementit}

% proof: proof/proposition-multiplication-n
\begin{proof}[{Proof}]\phantomsection\label{bp35}%
We use the isomorphism \([n]^*\mathcal N\simeq\mathcal N^{\otimes n^2}\) twice: on curves in the fibres of \([n]\) it shows that \([n]\) is finite, and in Euler characteristics it computes the degree. In between, translation invariance shows that all fibres have the same length, which makes \([n]\) finite locally free. Finally, whether \([n]\) is étale is decided by its differential.

Since finiteness, flatness and the degree may be checked after base change to an algebraic closure, we may assume that \(k\) is algebraically closed. Choose a symmetric ample \(\mathcal N\) (Lemma \hyperref[bp29]{B.10}); then \([n]^*\mathcal N\simeq\mathcal N^{\otimes n^2}\) by Proposition \hyperref[bp31]{B.11}. If a fibre of \([n]\) had positive dimension, it would contain a proper integral curve. Since \([n]\) maps that curve to a point, the restriction of \([n]^*\mathcal N\) to it is trivial, whereas the restriction of \(\mathcal N^{\otimes n^2}\) has positive degree. Hence every fibre has dimension zero, and since \([n]\) is proper, it is finite. Its image is a closed subset of dimension \(g\), hence equal to the integral variety \(A\) (Lemma \hyperref[bp26]{B.9}). After base change to an algebraic closure of its residue field, every fibre of \([n]\) becomes a translate of \(A[n]\). Indeed, if \(y=nx_0\), then translation by \(x_0\) carries \(A[n]=[n]^{-1}(0)\) isomorphically onto \([n]^{-1}(y)\), because \([n]\) is a homomorphism. Hence all fibres have the same degree \(r=\dim_k\Gamma(A[n],\mathcal O_{A[n]})\).

To get flatness from the constancy of the fibre degree, we use the following general fact. A finite morphism \(f\) to a reduced Noetherian scheme whose fibres all have degree \(r\) is finite locally free of rank \(r\). Indeed, near any point Nakayama's lemma gives a surjection \(\mathcal O^{\oplus r}\to f_*\mathcal O\), since the fibre at that point has dimension \(r\). This surjection is an isomorphism at each generic point, where the local ring is a field and both sides are vector spaces of dimension \(r\). Hence on an affine open every coordinate of a section of its kernel lies in all minimal primes, and is therefore zero because the ring is reduced. Since \(A\) is smooth, hence reduced, \([n]\) is finite locally free.

Recall that for a locally free sheaf \(\mathcal E\) of rank \(r\) on \(A\), the function \(m\mapsto\chi(A,\mathcal E\otimes\mathcal N^{\otimes m})\) is a polynomial of degree \(g\) with leading coefficient \(r\deg_{\mathcal N}(A)/g!\). Here \(\deg_{\mathcal N}(A)=(\mathcal N^g)\) is the degree of \(A\) with respect to \(\mathcal N\), a positive integer because \(\mathcal N\) is ample. We compare the coefficients of \(m^g\) in two expressions for \(\chi(A,[n]^*\mathcal N^{\otimes m})\). Since \([n]\) is finite, this Euler characteristic equals \(\chi(A,[n]_*[n]^*\mathcal N^{\otimes m})\), which by the projection formula is \(\chi(A,\mathcal N^{\otimes m}\otimes[n]_*\mathcal O_A)\). As \([n]_*\mathcal O_A\) is locally free of rank \(\deg[n]\), the coefficient is \(\deg[n]\deg_{\mathcal N}(A)/g!\). Since \([n]^*\mathcal N^{\otimes m}\simeq\mathcal N^{\otimes n^2m}\), it is also \(n^{2g}\deg_{\mathcal N}(A)/g!\). Hence

\[
\deg[n]\,\deg_{\mathcal N}(A)
=n^{2g}\deg_{\mathcal N}(A).
\]

Since \(\deg_{\mathcal N}(A)>0\), we conclude that \(\deg[n]=n^{2g}\); compare [Stacks, \href{https://stacks.math.columbia.edu/tag/0BFG}{0BFG}]. It remains to treat étaleness. The differential of the group law at \((0,0)\) is addition of tangent vectors. By induction, the differential of \([n]\) at the identity is \(n\operatorname{id}\). Since \([n]\) commutes with translations, the same holds at every point. If \(n\) is invertible, \([n]\) is therefore étale. Base change along the identity section gives the assertion for \(A[n]\).
\end{proof}

% text: prose-B-11
\phantomsection\label{bp36}%
It follows that the group \(A(\bar k)\) is divisible. Indeed, since \([n]\) is finite and surjective, its fibre over a point of \(A(\bar k)\) is a nonempty finite \(\bar k\)-scheme, and so it has a \(\bar k\)-point.


% object: sec-Tate-modules (blueprint/chap-Abelian-varieties/sec-Tate-modules/_section.md), content hidden
\subsection*{\texorpdfstring{B.3\quad Tate modules}{B.3 Tate modules}}
\phantomsection\label{bp37}%
\addcontentsline{toc}{subsection}{\texorpdfstring{B.3\quad Tate modules}{B.3 Tate modules}}
% text: prose-B-12
\phantomsection\label{bp38}%
Fix a prime \(\ell\) invertible in \(k\). With this prime, compatible systems of \(\ell\)-power torsion points of \(A\) form the Tate module \(T_\ell A\), a \(\mathbf Z_\ell\)-module with an action of \(\operatorname{Gal}(\bar k/k)\). For example, for a complex torus \(\mathbf C^g/\Lambda\) it is \(\Lambda\otimes\mathbf Z_\ell\), the first homology with \(\ell\)-adic coefficients. In general, Proposition \hyperref[bp40]{B.14} shows that \(T_\ell A\) is free of rank \(2g\) over \(\mathbf Z_\ell\). The input is the orders of the groups \(A[\ell^n](\bar k)\), which are known from Proposition \hyperref[bp34]{B.12}, and these orders force each of these groups to be free over \(\mathbf Z/\ell^n\) of rank \(2g\). Moreover, the transition maps are surjective because \(A(\bar k)\) is divisible.

Although Frobenius can have many fixed points on the finite levels, no power of it has a nonzero fixed vector in the Tate module (Proposition \hyperref[bp43]{B.15}). For example, for the elliptic curve \(E_0:y^2=x^3-x\) over \(\mathbf F_3\), the three roots of \(x^3-x\) lie in \(\mathbf F_3\). So all of \(E[2]\) is rational, and Frobenius acts trivially on it. On \(T_2E\), however, a fixed vector would be an unbounded compatible system of \(2\)-power torsion points in the finite group \(E_0(\mathbf F_{3^m})\). For the Jacobian of a curve, Proposition \bpout{14.3} turns this into the statement that Frobenius has no fixed class in \(H^1\) of the curve.

% object: definition-tate-module (blueprint/chap-Abelian-varieties/sec-Tate-modules/definition-tate-module/_section.md)
\phantomsection\label{bp39}%
\begin{bpstatement}{Definition~B.13}{ (Tate module)}
Let \(A\) be an abelian variety over \(k\) and let \(\ell\) be a prime invertible in \(k\). The \emph{Tate module} of \(A\) is

\[
T_\ell A=\varprojlim_n A[\ell^n](\bar k),
\]

the limit taken along the maps \([\ell]:A[\ell^{n+1}](\bar k)\to
A[\ell^n](\bar k)\), and we put \(V_\ell A=T_\ell A\otimes_{\mathbf Z_\ell}\mathbf Q_\ell\).
\end{bpstatement}

% object: proposition-integral-tate-module (blueprint/chap-Abelian-varieties/sec-Tate-modules/proposition-integral-tate-module/_section.md)
\phantomsection\label{bp40}%
\begin{bpstatementit}{Proposition~B.14}{ (The integral Tate module)}
Let \(A\) be an abelian variety of dimension \(g\) over \(k\), and let \(\ell\) be a prime invertible in \(k\). There are noncanonical isomorphisms

\[
A[\ell^n](\bar k)\simeq(\mathbf Z/\ell^n)^{2g},\qquad
T_\ell A\simeq\mathbf Z_\ell^{2g}.
\]

The canonical projection induces an isomorphism \(T_\ell A/\ell^nT_\ell A\simeq A[\ell^n](\bar k)\). The groups \(A[\ell^n](\bar k)\), the transition maps and the projections from \(T_\ell A\) are \(\operatorname{Gal}(\bar k/k)\)-equivariant, and they do not change when \(\bar k\) is replaced by a larger algebraically closed field.
\end{bpstatementit}

% proof: proof/proposition-integral-tate-module
\begin{proof}[{Proof}]\phantomsection\label{bp41}%
A finite abelian group of order \(\ell^{2gn}\) killed by \(\ell^n\) need not be free over \(\mathbf Z/\ell^n\); for \(g\geq1\) the group \((\mathbf Z/\ell^2)^{2g-2}\oplus(\mathbf Z/\ell)^{4}\) has order \(\ell^{4g}\) and is killed by \(\ell^2\). For \(A[\ell^n]\), however, this cannot happen because we know the orders at every level together with the exact sequences given by \([\ell]\). Indeed, we use them to lift a basis of the \(\ell\)-torsion step by step and check by counting that the lifts form a basis at each level.

\emph{The \(\ell\)-torsion and the exact sequences.} The group \(A[\ell](\bar k)\) is a vector space of dimension \(2g\) over \(\Fell\), since it is killed by \(\ell\) and has order \(\ell^{2g}\) by Proposition \hyperref[bp34]{B.12}. Moreover there are exact sequences

\[
0\longrightarrow A[\ell](\bar k)
\longrightarrow A[\ell^{n+1}](\bar k)
\xrightarrow{[\ell]}A[\ell^n](\bar k)\longrightarrow0,
\]

whose surjectivity holds because \(A(\bar k)\) is divisible: every point of \(A[\ell^n](\bar k)\) is \(\ell\) times a point, which then lies in \(A[\ell^{n+1}](\bar k)\).

\emph{Lifting bases.} Choose a basis of \(A[\ell](\bar k)\) and lift it successively along \([\ell]\); the lifts form a basis at each level. Suppose that \(x_1,\ldots,x_{2g}\in A[\ell^{n+1}](\bar k)\) lift a basis of \(A[\ell^n](\bar k)\) obtained in this way, so that \(\ell^nx_1,\ldots,\ell^nx_{2g}\) is the chosen basis of \(A[\ell](\bar k)\). To see that the \(x_i\) are independent, note first that all coefficients of a relation between them are divisible by \(\ell^n\), as we see by applying \([\ell]\) to it. The relation then becomes a relation between the elements \(\ell^nx_i\) of the chosen basis of \(A[\ell](\bar k)\), so the coefficients are even divisible by \(\ell^{n+1}\). As \(A[\ell^{n+1}](\bar k)\) has order \(\ell^{2g(n+1)}\), the \(x_i\) therefore form a basis of it over \(\mathbf Z/\ell^{n+1}\). Passing to the inverse limit of these compatible bases, we obtain both isomorphisms and the description of \(T_\ell A/\ell^nT_\ell A\).

\emph{Equivariance and independence of \(\bar k\).} The groups and maps in the statement commute with the action of \(\operatorname{Gal}(\bar k/k)\), because the group schemes \(A[\ell^n]\) and the maps \([\ell]\) are defined over \(k\). In particular the projection \(T_\ell A/\ell^nT_\ell A\to A[\ell^n](\bar k)\) is equivariant. The chosen bases need not be equivariant, which is why the isomorphisms with \((\mathbf Z/\ell^n)^{2g}\) and \(\mathbf Z_\ell^{2g}\) are noncanonical. Finally, the points of \(A[\ell^n]\) and the multiplication maps do not change under a further algebraically closed extension, because \(A[\ell^n]\) is finite étale and hence a finite disjoint union of points over \(\bar k\).
\end{proof}

% text: prose-B-13
\phantomsection\label{bp42}%
Over a finite field, Frobenius acts on \(T_\ell A\), and by the last proposition this action is compatible with its action on each \(A[\ell^n](\bar k)\). The next proposition shows that no power of Frobenius has a nonzero fixed vector. For the Jacobian of a curve, this is the input to Proposition \bpout{14.3}.

% object: proposition-no-frobenius-fixed-vector (blueprint/chap-Abelian-varieties/sec-Tate-modules/proposition-no-frobenius-fixed-vector/_section.md)
\phantomsection\label{bp43}%
\begin{bpstatementit}{Proposition~B.15}{ (No Frobenius-fixed vector)}
Let \(A_0/\mathbf F_q\) be an abelian variety, let \(A=A_0\times_{\mathbf F_q}\overline{\mathbf F}_q\), and let \(F\) be geometric Frobenius on \(T_\ell A\). For every \(m>0\),

\[
(T_\ell A)^{F^m}=0,
\qquad (V_\ell A)^{F^m}=0.
\]

Consequently \(F^m-1\) is an automorphism of the finite-dimensional \(\mathbf Q_\ell\)-vector space \(V_\ell A\), and hence of \(V_\ell A\otimes_{\mathbf Q_\ell}E\) for every field extension \(E\) of \(\mathbf Q_\ell\).
\end{bpstatementit}

% proof: proof/proposition-no-frobenius-fixed-vector
\begin{proof}[{Proof}]\phantomsection\label{bp44}%
A fixed vector is a compatible system of torsion points, each rational over \(\mathbf F_{q^m}\). But the rational points form a finite group, so the orders of the \(\ell\)-power torsion points in it are bounded, and compatibility under \([\ell]\) then forces the system to vanish. The statement for \(V_\ell A\) then follows by clearing denominators.

Since an automorphism and its inverse have the same fixed vectors, we may use arithmetic Frobenius instead. On \(A(\overline{\mathbf F}_q)\) its \(m\)th power fixes exactly \(A_0(\mathbf F_{q^m})\). This group is finite, because on each affine open of a finite affine cover a rational point is determined by finitely many coordinates in the finite field. Hence the \(\ell\)-primary part of \(A_0(\mathbf F_{q^m})\) is killed by some \(\ell^e\). Let \((t_n)\) be a fixed compatible sequence. Each \(t_{n+e}\) lies in this \(\ell\)-primary part, and since the transition maps are \([\ell]\), we find

\[
t_n=\ell^e t_{n+e}=0
\]

for every \(n\). This proves that \((T_\ell A)^{F^m}=0\).

Since \(T_\ell A\) is free by Proposition \hyperref[bp40]{B.14}, it injects into \(V_\ell A\). A suitable power of \(\ell\) times a fixed vector of \(V_\ell A\) is a fixed vector of \(T_\ell A\), and so it vanishes. Hence \((V_\ell A)^{F^m}=0\). In other words, \(F^m-1\) is injective on the finite-dimensional space \(V_\ell A\), and therefore bijective.
\end{proof}



\clearpage\phantomsection\hypertarget{bp-outside}{}%
\section*{Outside this export}
\addcontentsline{toc}{section}{Outside this export}
References marked \ensuremath{{}^\ast} are to parts of the blueprint this document leaves out:
\begin{description}
\item[Proposition 14.3] No fixed class in degree one, in Ⅲ Weights and the Riemann hypothesis (\bpslug{proposition-{}\allowbreak{}no-{}\allowbreak{}fixed-{}\allowbreak{}class-{}\allowbreak{}degree-{}\allowbreak{}one})
\item[Lemma A.4] Uniqueness of extensions from a dense open, in Appendices (\bpslug{lemma-{}\allowbreak{}uniqueness-{}\allowbreak{}maps-{}\allowbreak{}dense-{}\allowbreak{}open})
\end{description}
\clearpage\phantomsection%
\section*{References}
\addcontentsline{toc}{section}{References}
\bpbibitem{BBD82}{bp46}{A. Beilinson, J. Bernstein, and P. Deligne. \emph{Faisceaux pervers}, volume 100 of \emph{Astérisque}. 1982.}

\bpbibitem{BS15}{bp47}{B. Bhatt and P. Scholze. The pro-étale topology for schemes. \emph{Astérisque}, 369:99–201, 2015. \href{https://people.mpim-bonn.mpg.de/scholze/proetale.pdf}{Author's manuscript}.}

\bpbibitem{Con07}{bp48}{B. Conrad. Deligne's notes on Nagata compactifications. \emph{Journal of the Ramanujan Mathematical Society}, 22:205–257, 2007. \href{https://math.stanford.edu/\~conrad/papers/nagatafinal.pdf}{Author's manuscript}.}

\bpbibitem{Del74}{bp49}{P. Deligne. La conjecture de Weil. I. \emph{Publ. Math. IHÉS}, 43:273–307, 1974.}

\bpbibitem{Del80}{bp50}{P. Deligne. La conjecture de Weil. II. \emph{Publ. Math. IHÉS}, 52:137–252, 1980.}

\bpbibitem{Den22}{bp51}{C. Deninger. There is no “Weil-”cohomology theory with real coefficients for arithmetic curves. \emph{Ann. Sc. Norm. Super. Pisa Cl. Sci.}, 2022. To appear. arXiv:\href{https://arxiv.org/abs/2204.02714}{2204.02714}.}

\bpbibitem{Deu41}{bp52}{M. Deuring. Die Typen der Multiplikatorenringe elliptischer Funktionenkörper. \emph{Abh. Math. Sem. Univ. Hamburg}, 14:197–272, 1941.}

\bpbibitem{Eke90}{bp53}{T. Ekedahl. On the adic formalism. In \emph{The Grothendieck Festschrift, Vol. II}, volume 87 of \emph{Progress in Mathematics}, pages 197–218. Birkhäuser, 1990.}

\bpbibitem{Has36}{bp54}{H. Hasse. Zur Theorie der abstrakten elliptischen Funktionenkörper. I, II, III. \emph{J. Reine Angew. Math.}, 175:55–62, 69–88, 193–208, 1936.}

\bpbibitem{Hum72}{bp55}{J. E. Humphreys. \emph{Introduction to Lie Algebras and Representation Theory}, volume 9 of \emph{Graduate Texts in Mathematics}. Springer, 1972.}

\bpbibitem{Ill94}{bp56}{L. Illusie. Autour du théorème de monodromie locale. \emph{Astérisque}, 223:9–57, 1994. §4. \href{https://www.numdam.org/item/AST_1994__223__9_0/}{Original article}.}

\bpbibitem{Kat88}{bp57}{N. M. Katz. Travaux de Laumon. \emph{Astérisque}, 161–162:105–132, 1988. Séminaire Bourbaki, exp. 691.}

\bpbibitem{KW01}{bp58}{R. Kiehl and R. Weissauer. \emph{Weil conjectures, perverse sheaves and \(\ell\)-adic Fourier transform}, volume 42 of \emph{Ergebnisse der Mathematik}. Springer, 2001.}

\bpbibitem{Lau87}{bp59}{G. Laumon. Transformation de Fourier, constantes d'équations fonctionnelles et conjecture de Weil. \emph{Publ. Math. IHÉS}, 65:131–210, 1987.}

\bpbibitem{Mil86a}{bp45}{J. S. Milne. \href{https://www.jmilne.org/math/xnotes/AVs.pdf}{Abelian Varieties}. In \emph{Arithmetic Geometry}. 1986. Corrected online version.}

\bpbibitem{Mil86b}{bp60}{J. S. Milne. \href{https://www.jmilne.org/math/xnotes/JVs.pdf}{Jacobian Varieties}. In \emph{Arithmetic Geometry}. 1986. Corrected online version.}

\bpbibitem{Ran39}{bp61}{R. A. Rankin. Contributions to the theory of Ramanujan's function \(\tau(n)\) and similar arithmetical functions. II. \emph{Proc. Cambridge Philos. Soc.}, 35:357–372, 1939.}

\bpbibitem{SGA4}{bp62}{M. Artin, A. Grothendieck, and J.-L. Verdier, editors. \emph{Théorie des topos et cohomologie étale des schémas}, volumes 269, 270, 305 of \emph{Lecture Notes in Mathematics}. Springer, 1972–1973. SGA 4.}

\bpbibitem{SGA4½}{bp63}{P. Deligne. \emph{Cohomologie étale}, volume 569 of \emph{Lecture Notes in Mathematics}. Springer, 1977. SGA 4½; \emph{Théorèmes de finitude}, §3, and \emph{Rapport sur la formule des traces}, §3, and \emph{Sommes trigonométriques}, §1.}

\bpbibitem{SGA5}{bp64}{A. Grothendieck, editor. \emph{Cohomologie \(\ell\)-adique et fonctions \(L\)}, volume 589 of \emph{Lecture Notes in Mathematics}. Springer, 1977. SGA 5.}

\bpbibitem{SGA7II}{bp65}{P. Deligne and N. Katz, editors. \emph{Groupes de monodromie en géométrie algébrique}, volume 340 of \emph{Lecture Notes in Mathematics}. Springer, 1973. SGA 7 II; Exposé XIII, (2.1.7).}

\bpbibitem{Stacks}{bp66}{The Stacks Project Authors. \href{https://stacks.math.columbia.edu}{The Stacks Project}. Individual results are cited by their stable tags in the text.}

\bpbibitem{Var05}{bp67}{Y. Varshavsky. Lefschetz–Verdier trace formula and a generalization of a theorem of Fujiwara. 2005. Version 2. arXiv:\href{https://arxiv.org/abs/math/0505564}{math/0505564}. \href{https://arxiv.org/abs/math/0505564v2}{Author's preprint}.}

\bpbibitem{Wei48a}{bp68}{A. Weil. \emph{Sur les courbes algébriques et les variétés qui s'en déduisent}. Hermann, Paris, 1948.}

\bpbibitem{Wei48b}{bp69}{A. Weil. \emph{Variétés abéliennes et courbes algébriques}. Hermann, Paris, 1948.}

\bpbibitem{Wei49}{bp70}{A. Weil. Numbers of solutions of equations in finite fields. \emph{Bull. Amer. Math. Soc.}, 55:497–508, 1949.}

\end{document}
